\documentclass[12pt]{article}
\usepackage[utf8]{inputenc}
\usepackage[T1]{fontenc}
\usepackage[margin=1in]{geometry}
\usepackage{setspace}
\usepackage{xcolor}
\usepackage[colorlinks=true, urlcolor=blue!60!black, linkcolor=blue!60!black]{hyperref}
\setlength{\parindent}{2em}
\setlength{\parskip}{1.25em}
\renewcommand{\baselinestretch}{1.0}
\begin{document}
\begin{center}
{\LARGE \textit{Clapper v. Amnesty International USA}, 568 U.S. 398 (2013)\par}
\vspace{0.8em}
{\large Case Brief\par}
\end{center}
\vspace{1em}

\section{Memory Jogger}

In \href{https://supreme.justia.com/cases/federal/us/568/398/}{\textit{Clapper v. Amnesty International USA}, 568 U.S. 398 (2013)}, a 5--4 Supreme Court held that American attorneys, journalists, and human rights organizations lacked Article III standing to challenge Section 702 of the Foreign Intelligence Surveillance Act (50 U.S.C. \S~1881a) because their fear that the government would intercept their overseas communications rested on a five-link ``speculative chain of possibilities'' rather than ``certainly impending'' harm, and they could not ``manufacture standing'' by incurring costly protective measures against non-imminent surveillance.

\section{Detailed Case Facts}

In 1978, Congress enacted the Foreign Intelligence Surveillance Act (FISA), 50 U.S.C. \S~1801 \textit{et seq.}, creating the Foreign Intelligence Surveillance Court (FISC) to authorize electronic surveillance for foreign intelligence purposes upon a showing of probable cause that the target was a foreign power or an agent of a foreign power and that the specific facilities to be monitored were being used by that target (\href{https://tile.loc.gov/storage-services/service/ll/usrep/usrep568/usrep568398/usrep568398.pdf}{568 U.S. at 402--403}). Following the September 11, 2001 attacks, the Executive Branch conducted warrantless surveillance of international communications involving suspected members or agents of al Qaeda; after a FISC judge narrowed judicial authorization for that program in 2007, Congress enacted the FISA Amendments Act of 2008 (FAA), Pub. L. No. 110-261, 122 Stat. 2436, adding Section 702, codified at 50 U.S.C. \S~1881a (\href{https://supreme.justia.com/cases/federal/us/568/398/}{Justia Opinion, Part I.A}).

Section 1881a empowered the Attorney General and the Director of National Intelligence to jointly authorize, for up to one year, the targeting of non-U.S. persons reasonably believed to be located outside the United States to acquire foreign intelligence information. Unlike traditional FISA, \S~1881a eliminated the requirements that the government demonstrate probable cause that each target is a foreign power or agent of a foreign power or identify the specific facilities to be monitored, instead requiring FISC approval of programmatic ``targeting'' and ``minimization'' procedures and compliance with the Fourth Amendment (\href{https://supreme.justia.com/cases/federal/us/568/398/}{50 U.S.C. \S\S~1881a(a), (b)(5), (g)(2), (i)(2)--(3)}). Critically, \S~1881a barred the government from intentionally targeting any person known to be in the United States or any U.S. person reasonably believed to be located abroad (\S\S~1881a(b)(1)--(3)).

Respondents were United States persons---attorneys and human rights, labor, legal, and media organizations---whose work required sensitive and sometimes privileged telephone and e-mail communications with foreign clients, sources, witnesses, and colleagues abroad. Among the individual plaintiffs were attorney Scott McKay, who represented Sami Omar Al-Hussayen (whose prior communications had been intercepted roughly 10,000 times by phone and 20,000 times by e-mail under pre-2008 FISA authority) and Guant\'anamo Bay detainee Khalid Sheik Mohammed; attorney Sylvia Royce, who represented Guant\'anamo detainee Mohammedou Ould Salahi and communicated with his brother in Germany; journalist Christopher Hedges; and human rights researcher Joanne Mariner, who communicated with former detainees, lawyers, and activists in Jordan, Egypt, Pakistan, Afghanistan, and the Gaza Strip (\href{https://supreme.justia.com/cases/federal/us/568/398/}{Justia Opinion, Part I.B \& Dissent Part II.B}). Respondents alleged that because they communicated with likely targets of \S~1881a surveillance, the statute forced them to cease certain electronic communications, speak in generalities, and incur costly and burdensome measures---including planning foreign travel for in-person meetings---to protect confidentiality.

\section{Procedural History}

On July 10, 2008---the same day Congress enacted the FISA Amendments Act---respondents filed suit in the United States District Court for the Southern District of New York seeking a declaratory judgment that \S~1881a facially violated the Fourth Amendment, the First Amendment, Article III, and separation-of-powers principles, along with a permanent injunction against surveillance under \S~1881a (\href{https://clearinghouse.net/case/12972/}{Civil Rights Litigation Clearinghouse}).

On cross-motions for summary judgment, the District Court (Judge John G. Koeltl) granted summary judgment for the government, holding that respondents lacked Article III standing because they offered only an abstract, subjective fear of monitoring under \S~1881a and could not bootstrap that fear into standing through self-imposed protective costs (\textit{Amnesty Int'l USA v. McConnell}, 646 F. Supp. 2d 633, 635, 645--653 (S.D.N.Y. 2009); \href{https://clearinghouse.net/case/12972/}{Clearinghouse Summary}).

A panel of the United States Court of Appeals for the Second Circuit (opinion by Judge Gerard E. Lynch) reversed, holding that respondents established standing on two independent grounds: (1) an ``objectively reasonable likelihood'' that their communications would be intercepted under \S~1881a in the future, and (2) present economic and professional injuries incurred to avoid a reasonable fear of future unlawful surveillance (\textit{Amnesty Int'l USA v. Clapper}, 638 F.3d 118, 133--139 (2d Cir. 2011); \href{https://supreme.justia.com/cases/federal/us/568/398/}{Justia Opinion, Part I.C}). The Second Circuit denied rehearing en banc by an equally divided 6--6 vote, over dissents by Judge Reena Raggi and Chief Judge Dennis Jacobs (\textit{Amnesty Int'l USA v. Clapper}, 667 F.3d 163 (2d Cir. 2011)).

The Supreme Court granted certiorari on May 21, 2012 (\href{https://www.supremecourt.gov/docketfiles/11-1025.htm}{Supreme Court Docket No. 11-1025}), heard oral argument on October 29, 2012 (\href{https://www.supremecourt.gov/oral_arguments/argument_transcripts/2012/11-1025.pdf}{Oral Argument Transcript}), and issued its decision on February 26, 2013, reversing the Second Circuit and remanding for further proceedings (\href{https://www.oyez.org/cases/2012/11-1025}{Oyez}).

\section{Judicial Votes}

\textbf{Vote:} 5--4, reversing and remanding the judgment of the Second Circuit.

\textbf{Majority Opinion:} Justice Samuel A. Alito, Jr., joined by Chief Justice John G. Roberts, Jr., and Justices Antonin Scalia, Anthony M. Kennedy, and Clarence Thomas.

\textbf{Concurrences:} None.

\textbf{Dissenting Opinion:} Justice Stephen G. Breyer, joined by Justices Ruth Bader Ginsburg, Sonia Sotomayor, and Elena Kagan (\href{https://supreme.justia.com/cases/federal/us/568/398/}{Justia Syllabus}).

\section{Holding}

\textbf{Judgment:} Reversed and remanded. Respondents lacked Article III standing at the summary judgment stage to seek declaratory and injunctive relief against surveillance authorized by 50 U.S.C. \S~1881a.

\textbf{Rule Announced:} To establish Article III standing based on prospective harm, a plaintiff must show that the threatened injury is ``certainly impending'' (or, as noted in Footnote 5, poses a ``substantial risk'' of occurring), not merely an ``objectively reasonable likelihood'' resting on a highly attenuated chain of speculative possibilities involving independent decisionmakers. Furthermore, plaintiffs cannot manufacture standing by inflicting present economic or professional costs on themselves in response to fears of hypothetical future harm that is not certainly impending, nor are such self-inflicted costs fairly traceable to the challenged statute (\href{https://tile.loc.gov/storage-services/service/ll/usrep/usrep568/usrep568398/usrep568398.pdf}{568 U.S. at 409--416}).

\textbf{Holding vs. Dicta:} The Court's core holding is that respondents' five-step inferential chain failed to establish either an imminent injury in fact or fair traceability to \S~1881a at summary judgment, and that voluntary countermeasures taken to avoid a non-imminent threat do not confer standing. By contrast, two prominent passages operate as influential dicta or alternative observations: (1) Footnote 5's acknowledgment that prior cases ``do not uniformly require plaintiffs to demonstrate that it is literally certain that the harms they identify will come about'' and sometimes recognize a ``substantial risk'' standard---which the Court later elevated into a disjunctive test in \textit{Susan B. Anthony List v. Driehaus}, 573 U.S. 149, 158 (2014) (\href{https://www.scotusblog.com/2014/06/commentary-susan-b-anthony-list-clapper-footnote-5-and-the-state-of-article-iii-standing-doctrine/}{SCOTUSblog Commentary}); and (2) Part IV.B's assurances that \S~1881a would not be insulated from judicial review because the government must provide notice under 50 U.S.C. \S\S~1806(c) and 1881e(a) when using \S~1881a-derived evidence in prosecutions, and because electronic communications service providers may challenge directives before the FISC under \S~1881a(h)(4).

\section{Analysis of Opinions}

\subsection{Majority Opinion (Alito, J.)}

Justice Alito grounded the standing inquiry in Article III's separation-of-powers function, emphasizing that the inquiry is ``especially rigorous'' when reaching the merits would require the Judiciary to review actions of the political branches in intelligence gathering and foreign affairs (\href{https://tile.loc.gov/storage-services/service/ll/usrep/usrep568/usrep568398/usrep568398.pdf}{568 U.S. at 408--409}). Rejecting the Second Circuit's ``objectively reasonable likelihood'' standard as inconsistent with precedent requiring that threatened injury be ``certainly impending,'' Justice Alito dismantled respondents' future-injury theory into a five-link ``highly attenuated chain of possibilities'' (568 U.S. at 410):
\begin{enumerate}
    \item The government must decide to target the communications of non-U.S. persons with whom respondents communicate (since \S~1881a prohibits targeting U.S. persons directly);
    \item The government must choose to invoke \S~1881a rather than traditional FISA (\S~1805), Executive Order 12333, or foreign intelligence partnerships;
    \item The Article III judges of the FISC must approve the government's proposed targeting and minimization procedures under \S~1881a and the Fourth Amendment;
    \item The government must actually succeed in intercepting those foreign contacts' communications; and
    \item Respondents themselves must be parties to the specific communications incidentally acquired.
\end{enumerate}

The majority stressed that at the summary judgment stage under \textit{Lujan v. Defenders of Wildlife}, 504 U.S. 555, 561 (1992), plaintiffs must set forth ``specific facts'' rather than mere assumptions, and reiterated the Court's ``usual reluctance to endorse standing theories that rest on speculation about the decisions of independent actors'' such as FISC judges (\href{https://supreme.justia.com/cases/federal/us/568/398/}{568 U.S. at 411--414}). Even if interception were imminent, the second link defeated traceability because respondents could only guess whether the government would use \S~1881a rather than unchallenged surveillance authorities. In Footnote 5, Justice Alito acknowledged cases invoking a ``substantial risk'' standard, but concluded respondents failed even that test because of their attenuated chain of inferences (\href{https://www.scotusblog.com/2014/06/commentary-susan-b-anthony-list-clapper-footnote-5-and-the-state-of-article-iii-standing-doctrine/}{568 U.S. at 414 n.5}).

Turning to respondents' second theory---present economic and professional costs---Justice Alito held that ``respondents cannot manufacture standing merely by inflicting harm on themselves based on their fears of hypothetical future harm that is not certainly impending'' (\href{https://supreme.justia.com/cases/federal/us/568/398/}{568 U.S. at 416}). Quoting Judge Raggi's dissent below, he warned that the Second Circuit's approach would let any enterprising plaintiff, ``for the price of a plane ticket,'' lower the Article III bar to a mere ``nonparanoid fear.'' Moreover, because the government had already monitored clients like Al-Hussayen under pre-2008 FISA authority, respondents' precautions were not fairly traceable to \S~1881a. Under \textit{Laird v. Tatum}, 408 U.S. 1 (1972), allegations of a ``subjective `chill' '' from knowledge of government intelligence gathering do not confer standing. Justice Alito distinguished \textit{Friends of the Earth, Inc. v. Laidlaw Environmental Services (TOC), Inc.}, 528 U.S. 167 (2000) (concededly ongoing unlawful pollution), \textit{Meese v. Keene}, 481 U.S. 465 (1987) (plaintiff directly regulated by an existing ``political propaganda'' label), and \textit{Monsanto Co. v. Geertson Seed Farms}, 561 U.S. 139 (2010) (concrete evidence of nearby genetically engineered crops and bee pollination range). Finally, invoking \textit{Valley Forge Christian College v. Americans United for Separation of Church and State, Inc.}, 454 U.S. 464, 489 (1982), he rejected the argument that standing must exist simply because otherwise no one could sue, while adding that FISC review, criminal-prosecution notice requirements under \S\S~1806(c) and 1881e(a), and service-provider challenges under \S~1881a(h)(4) preserved avenues for judicial review.

\subsection{Concurrences}

No separate concurring opinions were filed.

\subsection{Dissenting Opinion (Breyer, J.)}

Justice Breyer, joined by Justices Ginsburg, Sotomayor, and Kagan, argued that interception of at least some of respondents' communications was not ``speculative,'' but rather ``as likely to take place as are most future events that commonsense inference and ordinary knowledge of human nature tell us will happen'' (\href{https://supreme.justia.com/cases/federal/us/568/398/}{Justia Opinion, Dissent}). Focusing on the concrete affidavits of attorneys Scott McKay and Sylvia Royce and researcher Joanne Mariner, Justice Breyer identified four record-based factors demonstrating a ``very high likelihood'' of interception under \S~1881a: (1) respondents regularly engaged in electronic communications with foreigners abroad containing ``foreign intelligence information'' that \S~1881a (unlike pre-2008 FISA) newly authorized the government to collect without individualized probable cause; (2) both respondents (out of professional duty) and the government (out of national security imperative) had powerful motives to conduct and monitor those communications; (3) the government's past conduct---including intercepting 10,000 phone calls and 20,000 e-mails of McKay's client---proved its appetite for such surveillance; and (4) the NSA's immense technological capacity, combined with a 99.9\% FISC approval rate in 2011 (1,674 of 1,676 applications approved, with two withdrawn), made interception virtually assured if the government was ``doing its job'' (\href{https://supreme.justia.com/cases/federal/us/568/398/}{Justia Opinion, Dissent Part III}).

Doctrinally, Justice Breyer contended that ``certainty is not, and never has been, the touchstone of standing'' (\href{https://firstamendment.mtsu.edu/article/clapper-v-amnesty-international/}{MTSU First Amendment Encyclopedia}). Surveying Supreme Court precedent---including \textit{Pennell v. City of San Jose}, 485 U.S. 1 (1988), \textit{Blum v. Yaretsky}, 457 U.S. 991 (1982), \textit{Davis v. FEC}, 554 U.S. 724 (2008), \textit{MedImmune, Inc. v. Genentech, Inc.}, 549 U.S. 118 (2007), \textit{Duke Power Co. v. Carolina Environmental Study Group, Inc.}, 438 U.S. 59 (1978), and \textit{Monsanto}---as well as traditional anticipatory breach, nuisance, and declaratory judgment actions, he maintained that Article III requires only a ``reasonable probability'' or ``high probability'' of future harm, and that present mitigation costs incurred in response to such a probability independently satisfy injury in fact.

\section{Examples: Future Applications}

\begin{enumerate}
    \item \textbf{Same-Side Application (No Standing --- Speculative Data Breach Mitigation):} A healthcare network suffers a ransomware attack, but forensic audits confirm the attackers encrypted servers without exfiltrating patient records, and no patient data has appeared on the dark web or been misused. A patient sues in federal court seeking injunctive relief and reimbursement for credit-monitoring subscriptions she purchased out of fear of identity theft. Under \textit{Clapper}, she lacks Article III standing: her fear of future identity theft rests on a speculative chain of third-party criminal choices, and purchasing credit monitoring cannot manufacture standing to avoid a non-imminent threat (\href{https://www.wakeforestlawreview.com/2026/10/deciphering-charlottes-web-aceituno-v-dhs-and-the-burden-of-proving-an-ongoing-practice/}{Wake Forest Law Review}).

    \item \textbf{Same-Side Application (No Standing --- Generalized Chilling Effect from Secret Police Software):} A state legislature authorizes local police departments to purchase facial-recognition software for use in public spaces, subject to judicial warrants and statutory privacy guidelines. A civil liberties journalist who attends public protests sues to enjoin the statute, alleging she now wears disguises and avoids protests because police might target protest organizers and incidentally scan her face. Under \textit{Clapper}, she lacks standing because she offers no specific facts that police are targeting her or her associates under the new statute rather than existing cameras, and her self-imposed precautions reflect a subjective ``chill'' insufficient under \textit{Laird} and \textit{Clapper}.

    \item \textbf{Opposite-Side Application (Standing --- Targeted Pre-Enforcement Speech Threat):} A state statute criminalizes making ``misleading'' statements about ballot initiatives and empowers a state commission to hold probable-cause hearings upon any private complaint. An advocacy group that was previously hauled before the commission for a similar ad plans to run an identical ad next month and has already received a warning letter from an opponent threatening a new complaint. Under \textit{Clapper}'s Footnote 5 as applied in \href{https://www.scotusblog.com/2014/06/commentary-susan-b-anthony-list-clapper-footnote-5-and-the-state-of-article-iii-standing-doctrine/}{\textit{Susan B. Anthony List v. Driehaus}, 573 U.S. 149 (2014)}, the group has standing because it is directly regulated by the statute and faces a ``substantial risk'' of burdensome enforcement proceedings backed by criminal sanctions.

    \item \textbf{Opposite-Side Application (Standing --- Confirmed Data Exfiltration and Active Misuse):} Hackers breach a payroll processor, exfiltrate unencrypted Social Security numbers and bank routing numbers, and begin filing fraudulent tax returns for 15\% of the affected employees within two weeks. An affected employee whose data was confirmed stolen---though not yet misused---sues for injunctive relief and mitigation expenses. Unlike in \textit{Clapper}, the unauthorized acquisition has already occurred and active criminal exploitation is underway, establishing a ``substantial risk'' of imminent identity theft that renders reasonable mitigation costs traceable to the breach (\href{https://www.law.cornell.edu/constitution-conan/article-3/section-2/clause-1/actual-or-imminent-injury}{Cornell LII Constitution Annotated}).

    \item \textbf{Fence-Sitter Application (Close Call --- Confirmed Bulk Surveillance Infrastructure Without Confirmed Individual Interception):} Following official government declassification, an international news organization learns that an intelligence agency conducts ``upstream'' collection on three of the four major transatlantic fiber-optic cables carrying its reporters' overseas traffic, and that the agency scans all packets traversing those cables for foreign-intelligence selectors. The organization cannot prove that any specific reporter's e-mail was retained, because targeting lists remain classified, yet mathematically its traffic traverses the tapped cables daily. A court could find standing by distinguishing \textit{Clapper} on the ground that the first four links of Alito's speculative chain are eliminated by confirmed programmatic facts (cf. \href{https://law.justia.com/cases/federal/appellate-courts/ca4/20-1191/20-1191-2021-09-15.html}{\textit{Wikimedia Foundation v. NSA}, 14 F.4th 276 (4th Cir. 2021)}), or deny standing if it views the final link---incidental acquisition of a particular plaintiff's communications---as still unproven without invading state secrets.
\end{enumerate}

\section{Critique}

\textbf{Originalist and Textualist Perspectives:} Supporters of \textit{Clapper} argue that rigorous standing rules faithfully enforce Article III's limitation of ``the judicial Power'' to ``Cases'' and ``Controversies,'' ensuring federal courts resolve concrete disputes between ``proper parties'' seeking ``proper relief'' rather than issuing advisory opinions on national security statutes (\href{https://harvardlawreview.org/print/vol-137/proper-parties-proper-relief/}{Harvard Law Review, Vol. 137}). Yet some originalist and textualist scholars critique modern ``injury-in-fact'' doctrine---and \textit{Clapper}'s hyper-stringent ``certainly impending'' formulation---as an ahistorical twentieth-century gloss on Article III that conflates constitutional jurisdiction with the merits of whether a plaintiff holds a legal right or cause of action, while selectively closing courthouse doors in ways untethered to Founding-era equity and declaratory practice.

\textbf{Progressive and Doctrinal Perspectives:} Progressive scholars and civil-liberties critics contend that \textit{Clapper} creates a Catch-22 that effectively immunizes secret government surveillance from constitutional review: plaintiffs cannot prove they are being secretly wiretapped without discovery, yet they cannot obtain discovery without first proving they are being wiretapped (\href{https://www.lawfaremedia.org/article/why-clapper-matters-future-programmatic-surveillance}{Lawfare}; \href{https://journals.law.unc.edu/firstamendmentlawreview/wp-content/uploads/sites/8/2023/08/dickinson-piece-1.pdf}{First Amendment Law Review}). Critics also emphasize two doctrinal vulnerabilities. First, the majority's assurance that criminal defendants would receive notice of \S~1881a-derived evidence proved factually inaccurate at the time of argument: reporting in 2013 revealed that federal prosecutors had been withholding such notice, prompting the Justice Department to reverse course and issue its first FAA notice in October 2013 (\href{https://www.scotusblog.com/2013/07/wednesday-round-up-192/}{SCOTUSblog}). Second, critics point to stark doctrinal inconsistency in the Court's treatment of ``chilled speech'' and probabilistic harm---demanding near-certainty in \textit{Clapper} while permitting standing on far more attenuated risks of enforcement in First Amendment and regulatory challenges such as \textit{Susan B. Anthony List} (\href{https://texaslawreview.org/chill/}{Texas Law Review}; \href{https://harvardlawreview.org/print/vol-131/standing-surveillance-and-technology-companies/}{Harvard Law Review, Vol. 131}).

\textbf{Underlying Competing Values:} At its core, \textit{Clapper} reflects a clash between two foundational constitutional commitments: judicial restraint and executive operational secrecy in foreign intelligence on one side, versus constitutional accountability and the protection of First and Fourth Amendment rights against covert dragnet surveillance on the other. Requiring proof of ``certainly impending'' interception protects intelligence priorities from fishing expeditions and preserves separation of powers, but it risks leaving systemic constitutional violations unreviewable whenever government conduct is designed to remain invisible.

\section{Key Quotations}

\begin{quote}
``Thus, we have repeatedly reiterated that `threatened injury must be \textit{certainly impending} to constitute injury in fact,' and that `[a]llegations of \textit{possible} future injury' are not sufficient.''
\end{quote}
\noindent--- Justice Alito, Majority Opinion, \href{https://tile.loc.gov/storage-services/service/ll/usrep/usrep568/usrep568398/usrep568398.pdf}{568 U.S. at 409} (quoting \textit{Whitmore v. Arkansas}, 495 U.S. 149, 158 (1990)).

\begin{quote}
``We decline to abandon our usual reluctance to endorse standing theories that rest on speculation about the decisions of independent actors.''
\end{quote}
\noindent--- Justice Alito, Majority Opinion, \href{https://supreme.justia.com/cases/federal/us/568/398/}{568 U.S. at 414}.

\begin{quote}
``Our cases do not uniformly require plaintiffs to demonstrate that it is literally certain that the harms they identify will come about. In some instances, we have found standing based on a `substantial risk' that the harm will occur, which may prompt plaintiffs to reasonably incur costs to mitigate or avoid that harm.''
\end{quote}
\noindent--- Justice Alito, Majority Opinion, \href{https://supreme.justia.com/cases/federal/us/568/398/}{568 U.S. at 414 n.5}.

\begin{quote}
``In other words, respondents cannot manufacture standing merely by inflicting harm on themselves based on their fears of hypothetical future harm that is not certainly impending.''
\end{quote}
\noindent--- Justice Alito, Majority Opinion, \href{https://supreme.justia.com/cases/federal/us/568/398/}{568 U.S. at 416}.

\begin{quote}
``In my view, this harm is not `speculative.' Indeed it is as likely to take place as are most future events that commonsense inference and ordinary knowledge of human nature tell us will happen.''
\end{quote}
\noindent--- Justice Breyer, Dissenting Opinion, \href{https://supreme.justia.com/cases/federal/us/568/398/}{568 U.S. at 422}.

\end{document}
